\section{Directional mean weights and local nonlinear observations}\label{sec:average}
Pointwise bounds on a derivative repair are stronger than the following
actual spatial-mean condition. Fix $d\geq2$, $N_c\geq2$, and distinct
directions $i,j$. For nonnegative componentwise weights $w_n(k,(x,l))$, put
\[
 Q_n(A)=\ip{A}{K_{0,n}A}+\sum_{x,k,l}w_n(k,(x,l))
                  \|A_l(x+e_k)-A_l(x)\|^2,
 \qquad b_{ij}(n)=n^{-d}\sum_x w_n(i,(x,j)).
\]
The repair is the actual $G^*W_nG$ of the componentwise forward gradient.
In particular $Q_n\geq0$; signed weights are not included in this section.
Let $D_n=\ker\delta$ be the actual flat divergence-free space.

\begin{theorem}[Local nonlinear average-weight obstruction]\label{thm:average}
For each $n\geq1$, let $F_n$ be a finite-dimensional real normed space,
$C_n:E_n\to F_n$ be Fr\'echet differentiable at zero with $C_n(0)=0$,
and let $\Phi_n:F_n\to\R$ satisfy
\[
 \Phi_n(y)\leq K_n\|y\|^2\quad(\|y\|<\delta_n),\qquad \delta_n>0.
\]
Suppose one fixed $c>0$ and positive radii $r_n$ satisfy
\begin{equation}\label{eq:local-average}
 c\|A\|^2\leq Q_n(A)+\Phi_n(C_n(A))
       \quad(A\in D_n,\ \|A\|<r_n).
\end{equation}
For every integer $0\leq R<n$,
\begin{equation}\label{eq:mean-finite}
 c(2R-\dim F_n)\leq 2(1+b_{ij}(n))
       \sum_{s=1}^{R}|e^{2\pi\mathrm{i}s/n}-1|^2
 \leq \frac{8\pi^2(1+b_{ij}(n))R^3}{n^2}.
\end{equation}
Consequently, if $(1+b_{ij}(n))/n^2\to0$, then $\dim F_n\to\infty$.
No uniform bound on $K_n,\delta_n^{-1}$ or $r_n^{-1}$ is required.
\end{theorem}
\begin{proof}
Write $L_n=DC_n(0)$. If $v\in D_n\cap\ker L_n$, differentiability gives
$C_n(tv)=o(|t|)$. For all sufficiently small nonzero $t$, both the local
coercivity and penalty bound apply. Divide by $t^2$ and let $t\to0$ to get
\begin{equation}\label{eq:kernel-average}
 c\|v\|^2\leq Q_n(v).
\end{equation}
This argument is performed separately at each $n$; it does not impose
uniform neighborhoods or derivative remainder rates.

Use the two internal coordinates $u,v$ of Section 2 to regard a complex
number as a real internal vector. For $1\leq s\leq R$ and $a\in\{1,\mathrm{i}\}$,
define $\psi_{s,a}$ by taking only component $j$ nonzero and setting it to
$n^{-d/2}$ times the internal vector corresponding to
$a e^{2\pi\mathrm{i}s x_i/n}$. The finite-character orthogonality identity
and the real orthogonality of the two quadratures give an orthonormal
family of $2R$ fields. This also covers even $n$; the two quadratures use
two internal coordinates, not two allegedly independent scalar sine and
cosine functions. Since $i\ne j$, these fields are independent of $x_j$,
so $\delta\psi_{s,a}=0$. Their span $S_{R,n}$ is thus an actual
$2R$-dimensional subspace of $D_n$.

Every nonzero forward difference of $\psi_{s,a}$ is in direction $i$ and
has the same squared norm at every site, namely
$n^{-d}|e^{2\pi\mathrm{i}s/n}-1|^2$. The exact Hodge and repair energies are
therefore
\[
 Q_n(\psi_{s,a})=(1+b_{ij}(n))|e^{2\pi\mathrm{i}s/n}-1|^2.
\]
In particular the trace of the compression of $Q_n$ to $S_{R,n}$ is the
middle expression in \eqref{eq:mean-finite}; no invariance of $S_{R,n}$
under the weighted operator is asserted or needed.

For completeness, the positive trace/rank argument is as follows. Set
$U=S_{R,n}\cap\ker L_n$. Rank-nullity gives
$\dim U\geq2R-\dim F_n$. Extend an orthonormal basis of $U$ to one of
$S_{R,n}$. Equation \eqref{eq:kernel-average} bounds the terms on $U$
below by $c$; positivity of $Q_n$ makes every remaining diagonal term
nonnegative. Hence $c(2R-\dim F_n)\leq\operatorname{tr}_{S_{R,n}}Q_n$.
Finally $|e^{2\pi\mathrm{i}s/n}-1|\leq2\pi s/n\leq2\pi R/n$, which proves
the last inequality in \eqref{eq:mean-finite}.

Fix $R$. The right side tends to zero under the stated mean condition.
Eventually it is less than $c$. If $\dim F_n<2R$, integrality would make
the left side at least $c$. Thus $\dim F_n\geq2R$ eventually. Arbitrarily
large fixed $R$ prove the claimed divergence.
\end{proof}

The condition concerns one directional spatial average, not the maximum
of all weights. For example, a weight of size $n^{d+1}$ at one site in
component $(i,j)$ and zero elsewhere has $b_{ij}(n)=n$, while its maximum
is not subquadratic. Theorem~\ref{thm:average} still excludes bounded
output dimension under its other hypotheses. This example is a direct
calculation, not an additional formal endpoint. The nonnegative-weight
condition is essential to the trace argument as used here: the complement
of the derivative kernel must contribute nonnegative energy. This is not
a universal extension of older statements allowing signed weights with a
pointwise absolute cap.

\subsection{A constructed quadratic-weight boundary example}
The exponent $2$ cannot be replaced by an unrestricted allowance of
quadratic-size mean weights. Let
$\mu_j(A)=n^{-d}\sum_x A_j(x)$, with output in the same fixed space
$F=(\R^{N_c^2-1})^d$ for every $n$. A finite path comparison on the actual
periodic lattice gives
\[
 \|A-\operatorname{lift}\mu(A)\|^2\leq d n^2\|GA\|^2,
 \qquad \|\operatorname{lift}\mu(A)\|^2=n^d\|\mu(A)\|^2.
\]
One way to obtain the first inequality is to write each difference of
site values as a path of at most $dn$ forward edges, apply finite
Cauchy--Schwarz, and average the pairs of sites; the displayed constant
is deliberately nonoptimal. The exact modular-torus reindexing and its
transport to the physical cochain norm are also proved in the companion
formal construction.
Using $\|x+y\|^2\leq2\|x\|^2+2\|y\|^2$, choose actual weights
$w_n(k,(x,l))=2dn^2$ and the quadratic penalty
$\Phi_n(y)=2n^d\|y\|^2$. Then
\[
 \|A\|^2\leq\ip{A}{G^*W_nGA}+\Phi_n(\mu(A))
                  \leq Q_n(A)+\Phi_n(\mu(A))
\]
on the entire space, hence also on $D_n$. The mean observation is an
actual continuous linear map with finite-dimensional output type
independent of $n$, not a postulated observation certificate. Here
$b_{ij}(n)=2dn^2$ and the subquadratic hypothesis fails. The example
allows $K_n=2n^d$ to vary, exactly as the obstruction does. It proves
neither an optimal boundary constant nor positive coercivity for an
interacting gauge operator. The path/mean mechanism is classical;
see the primary exposition in \cite{levin-peres}.

\subsection{Why the penalty assumption remains}
The $C^1$ counterexample after Theorem~\ref{thm:tangent} also applies to
the actual nonnegative $Q_n$ and to $D_n$. A second useful counterexample
retains a polynomial observation and a globally Lipschitz penalty:
\[
 C(A)=\|A\|^2\in\R,\qquad \Phi(y)=c|y|,
 \qquad DC(0)=0,\qquad \Phi(C(A))=c\|A\|^2.
\]
Thus output dimension one suffices even when $Q_n=0$. No finite local
quadratic upper constant for $\Phi$ exists, since $c|y|/y^2=c/|y|$.
Likewise, $C(A)=\|A\|^{3/2}$ and $\Phi(y)=c|y|^{4/3}$ are both $C^1$
with zero derivatives at zero, yet their composition is exactly quadratic.
These are complete analytic counterexamples, not new Lean claims.

There is a distinct analytic regularity tradeoff, not a hypothesis removal
for merely differentiable observations. If, on each fixed derivative-kernel
direction, $C(tv)=t^2H_v+o(t^2)$, and $\Phi(0)=0$, $D\Phi(0)=0$, then
$\Phi(C(tv))=o(t^2)$. Indeed $\|C(tv)\|/t^2$ is bounded and
$|\Phi(y)|\leq\epsilon\|y\|$ near zero, including at zero; multiply the
two estimates and let $\epsilon\downarrow0$. A $C^2$ observation and a
nonnegative differentiable penalty vanishing at zero suffice for this
step, because a differentiable local minimum has zero derivative.
No actual observation's $C^2$ regularity is proved here, and no formal
composition of this alternate route is claimed.

The strengthened mean-weight result is an output-rank divergence theorem.
Its fixed-$R$ argument does not assert a positive density of observations
per lattice volume. Fixed $c>0$, nonnegative weights and the stated
penalty condition remain explicit. In particular no interacting or
continuum Yang--Mills hypothesis, ultraviolet activity estimate, or
physical spectral gap is supplied by this section.
